\chapter{放射防护管理}

\section{防护管理组织}

组长（院长刘绍华）、分管副院长、郑志刚为专职放射防护负责人，多科室成员组成管理小组。

\section{现有管理制度}

已建立：岗位责任制、设备操作规程、源管理制度、个人防护用品管理、患者告知、质控方案、辐射监测、应急预案等；
\textbf{建议}：根据新进设备修订针对性操作规程、制度上墙公示。

\section{职业健康管理}


全部 5 名工作人员按期在岗复训、个人剂量季度送检、每 2 年职业健康体检，体检结论均可继续放射作业。

\subsection{人员管理要求}

\begin{enumerate}
\item \textbf{培训}：上岗前≥4d、在岗每 2 年复训，江西省线上平台统一培训考核；
\item \textbf{个人剂量}：终身存档个人剂量档案；
\item \textbf{职业体检}：上岗前、在岗 2 年 / 次、离岗前体检，体检档案永久保存。
\end{enumerate}

\section{放射卫生档案分类存档}

\begin{enumerate}
\item 项目档案（预评价、控评、批复）→公卫科；
\item 证照档案（放射诊疗许可证）→公卫科；
\item 人员（培训、体检、剂量、资质）→公卫科 + 放疗科分存；
\item 设备、质控、监测、防护用品台账→设备科、放疗科。
\end{enumerate}

\section{小结}

组织健全、制度体系完善、人员健康监护闭环管理；持续更新制度并上墙。
